\chapter{Outlook}
\label{ch:outlook}

Thirteen chapters is a long way from the preface's claim, which was that giving
up the graph as the fundamental object and keeping the tree brings the whole
apparatus back. This chapter checks that claim against what was actually
delivered, records two things the book kept discovering by accident, and then
says what is not done --- at more length than is usual, because several of the
gaps are specific enough to act on.

\section{One recursion, many models}
\label{sec:onerecursion}

Table~\ref{tab:through} puts the through-line in one place. Down the columns
rather than across the rows: the recursion never changes. What
changes is what the message carries and what is summed inside a complex, and
everything else --- the layers, the direction of arrival, the excess bracket, the
$L\times L$ determinant --- is common to every row.

\begin{table}[b]
\centering\small
\begin{adjustbox}{max width=\linewidth}
\begin{tabular}{llll}
\hline\hline
chapter & the message is & the interior sum is & and one gets\\
\hline
\ref{ch:percolation}--\ref{ch:giant} & one number
  & which members are reachable & $\theta$, $\Lambda$, $B$\\
\ref{ch:epidemics} & one number
  & a within-group outbreak & $R^{*}$, the final size\\
\ref{ch:potts} & one number
  & a Potts partition ratio & $\tau_{c}(q,v)$\\
\ref{ch:statmech} & a measure
  & any $-\beta H$, over $2^{c}$ states & $\det(I-B)=0$\\
\ref{ch:ising} & a cavity field
  & all pairs, or unanimity & $T_{c}$, the AT line\\
\ref{ch:hittingset} & an integer, then a field
  & one forbidden configuration & $\ave{k}(c-1)=e$, $\rho$\\
\ref{ch:cover} & three states
  & the leaf-removal rule & $\ave{k}=e$, the core\\
\ref{ch:colouring} & a distribution over $q$
  & a hard colour constraint & $\ave{k}=(q-1)^{2}$\\
\ref{ch:satisfiability} & a warning, at two levels
  & one forbidden assignment & $\alpha=1$ at $k=2$\\
\ref{ch:overlap} & a function on a region
  & a region-graph enumeration & what the rest costs\\
\hline\hline
\end{tabular}
\end{adjustbox}
\caption{\textbf{One recursion.} Every chapter of Parts~II and~III is
Eqs.~\eqref{eq:up} and~\eqref{eq:emitted} with two things substituted. The
column that does not appear is the recursion itself, because it is the same in
every row.}
\label{tab:through}
\end{table}

Two rows deserve a second look, because they are the ones that make the table a
claim rather than a filing system.

Chapter~\ref{ch:potts} showed that the first and fifth rows are not merely
similar. The $q\to1$ limit of the Potts interior sum \emph{is} the reachability
generating function of Chapter~\ref{ch:percolation}, and $q=2$ \emph{is} the
Ising cavity field of Chapter~\ref{ch:ising}: one function, $\tau_{c}(q,v)$,
evaluated at two points. So the two halves of the book are not two subjects
sharing a method. Percolation came first because at $q=1$ the message is a
probability and probabilities multiply, which is arithmetic a reader already
has --- and for no deeper reason than that.

Chapter~\ref{ch:satisfiability}'s row is the only one where a complex contains
another complex, which is the freedom the definition was built for and the one
the rest of the book barely uses. Section~\ref{sec:flatteningcost} measures it: a
nested clause has a treelike three-level representation, its CNF flattening does
not, and a flat treatment misplaces the transition by seven to twenty-two per
cent. The deficit is manufactured by the \emph{encoding} rather than found in
the data, which is Chapter~\ref{ch:giant}'s moral in its strongest form --- the
tree calculation applied to the wrong object.

And Chapter~\ref{ch:overlap}'s row is the one that breaks the pattern
deliberately. Its message is a function on a region rather than a distribution
over a class, which is why that chapter is instance-level while every other row
is an ensemble. That single mismatch is the book's largest open problem, and
Sec.~\ref{sec:notdone} states it as such.

\section{Two things the book kept discovering}
\label{sec:threads}

Neither of these was planned. Each turned up in five independent places, which
is the only reason to trust them.

\paragraph{Name what is held fixed.} A comparison between a clustered structure
and an unclustered one is not a statement about clustering until one says what
the two have in common. The book found this five times, and the sign of the
effect changed twice.

Chapter~\ref{ch:data} found that a degree-matched control can carry \emph{more}
clique overlap than the real network, because heavy tails manufacture cliques at
hubs. Section~\ref{sec:helporhurt} found that triangles help percolation at
fixed link count and hurt it at fixed degree --- the link-matched comparison
lets the degree distribution broaden, and $q_{c}$ moves from $1/3$ to $0.4030$
when that is held down. Section~\ref{sec:joint} found three joint distributions
with identical marginals, hence identical threshold tensors, and thresholds
spread over sixty per cent. Section~\ref{sec:epidemics-lessons} found that
households help an epidemic when added on top of an unchanged global network and
hurt it at a fixed contact budget, $T_{c}$ going from $0.2500$ to $0.2929$ the
wrong way. And Sec.~\ref{sec:clustering} found that a Poisson link layer matched
only on \emph{mean} degree reverses the sign of the Ising result, because the
triangle construction carries excess degree $\ave{\bar k}=n+1$ rather than $n$
and a Poisson triangle layer loses none of its branches.

Five settings, one lesson, and it is not a lesson about chygraphs. It is a
lesson about null models, and the reason it recurs here is that promoting motifs
to complexes makes it very easy to build a control that differs in more than one
thing at once.

\paragraph{A linear instability is not always the transition.} The branching
condition $\det(I-B)=0$ locates the loss of stability of a trivial fixed point.
When the transition is continuous that is the transition; when it is not, it is
a spinodal, and quoting it is an error of definite sign and often of large size.

Section~\ref{sec:breakdown} met this first: interdependent networks acquire a
giant component by a saddle-node bifurcation while $Q\equiv1$ is still stable,
so $\Lambda=0$ does not locate anything. Section~\ref{sec:groupcontagion} met
the same exclusion for threshold contagion --- with $r\ge2$ a single seed
infects nobody, so the trivial fixed point is stable at every coupling.
Section~\ref{sec:transmission} met it for the Potts model above $q=2$, where the
branching condition sits $5.1$ to $22.8$ per cent above the transition at
$q=3$ to $6$. Section~\ref{sec:unanimity} met it where it does real damage, in the
unanimity interaction, and there the free energy built in
Sec.~\ref{sec:freeenergy} finally had to be used. And
Sec.~\ref{sec:twosat} is the extreme of the same statement: above $k=2$ the
warning map's derivative at the trivial fixed point is exactly zero at every
clause density, so there is no linear instability anywhere on the axis to
mistake for a transition.

The practical corollary is easy to get wrong in the direction that flatters the
calculation: critical slowing down mimics a coexistence window. An ordered branch that survives a finite iteration
budget just above the spinodal is not evidence of a first-order transition. The
test needs a free-energy crossing, and applying it moved the tricritical
boundary of Sec.~\ref{sec:unanimity} by one in cardinality.

\section{What the book establishes}
\label{sec:establishes}

Four results, stated as flatly as they can be.

\emph{Clustering lowers the Ising critical temperature at fixed degree, and a
triangle nevertheless transmits more than an edge.} Both halves are true;
Eq.~\eqref{eq:branch} carries a multiplicity as well as a transmission, and the
multiplicity wins. The effect is $-13.9\%$ at four neighbours, decays as
$1/\ave{k}$, and was confirmed outside the formalism by Wolff simulations at
$-14.2\%$.

\emph{The hard-field limit of covering problems is exact on graphs and wrong on
hypergraphs.} Disjoint $3$-hyperedges: no complex touches another, replica
symmetry cannot break, the density is $1/3$ by counting, and warning propagation
returns $1/2$. The cause is identified precisely --- the cavity field carries a
fraction $\mu/L$, which an integer ansatz cannot hold --- and the repair,
keeping the $O(1)$ fields, returns $1/3$ and reproduces the published closed
forms.

\emph{A complex of three or more nodes destroys the core-free branch.} The
polynomial condition for $\gamma=0$ factors as
$-(1-\delta)^{2}\sum_{j}(j+1)\delta^{j}$, identically zero at $c=2$ and strictly
negative above it. So above cardinality two leaf removal never certifies a
minimum cover --- which is exactly why the hard-field cover size stopped being
believable there, and the two statements turn out to be one.

\emph{Hyperbolic random graphs are predicted once the theory is applied to the
right object.} Not to the graph, whose clustering has defeated degree-based
mean-field theory for a decade, but to the chygraph of its maximal cliques, with
the ensemble measured rather than fitted. The degree-matched control has no core
anywhere; the chygraph has one everywhere, accounting for $77$ to $100$ per cent
of it.

That last result carries a qualification the book has stated twice and will
state a third time, because it is the one a reader is most likely to
over-remember. The \emph{existence} of the core follows from the algebra the
moment a single triangle enters the ensemble --- Sec.~\ref{sec:corefree}
guarantees it. Only the exponent could have failed, and there the evidence is
weaker than the central values suggest: the predicted exponent does not move
with the tail while the measured one does.

\section{What the book does not do}
\label{sec:notdone}

In rough order of how much it would take to close.

\emph{Overlapping complexes at the level of an ensemble.} This is the largest
gap and Chapter~\ref{ch:overlap} is entirely about it. Two partial answers exist
and neither closes it. The instance calculation is done and exact where the
clique structure is chordal; and Sec.~\ref{sec:merging} shows that merging the
offending complexes into one is exact whenever it terminates, with a percolation
criterion saying when it will --- which on the ensembles that motivated the
construction is only in the very sparse r\'egime. What is missing is a
message indexed by region \emph{type}, over a region graph that is itself
random. Two things would have to be supplied: the distribution of intersection
patterns among the maximal cliques of the ensemble, and what the parent-to-child
update becomes when averaged over it. And a warning is already in hand --- six
of twenty non-chordal instances do not converge at damping $0.999$, so an
ensemble treatment inherits an instability that damping alone will not fix.

\emph{Replica symmetry breaking, in general.} The book detects it and does not
do it.
Chapter~\ref{ch:hittingset} gives an entropy that proves the
replica-symmetric answer wrong where it goes negative, and says on which side of
the truth it then falls; Chapter~\ref{ch:ising} gives the de Almeida--Thouless
line; Chapter~\ref{ch:statmech} gives the sign test that says which solver
applies. What none of them gives is the one-step calculation itself, which for
the hitting set problem is exactly where the interesting régime begins.

\emph{Interdependent networks.} Section~\ref{sec:breakdown} showed that these
are outside the branching class: the giant component appears while the trivial
fixed point is still stable, so $\Lambda$ does not locate it. Mapping those
constructions to chygraphs, and identifying what replaces $\Lambda$ there, is
not done.

\emph{Threshold contagion, properly.} Section~\ref{sec:groupcontagion} gives the
reason the percolation machinery cannot carry a rule that needs several infected
members at once, and then solves a mean-field equation that keeps the mechanism.
What it does not give is a theory: that would need a message carrying more than
one number per inclusion, which is what Part~III is about, applied to a
dynamical rather than an equilibrium rule.

\emph{One-step replica symmetry breaking, where the subject begins.} Two
chapters end by saying this. Chapter~\ref{ch:colouring}'s threshold is a local
stability bound that for $q\ge4$ sits on the far side of the colourability
threshold; Chapter~\ref{ch:satisfiability} cannot even offer that, having no
symmetric fixed point to linearise about, and its hard-field substitute is blind
over a range that grows with $k$. Both would be closed by the same thing, and it
is the thing this book does not do.

\emph{The exponent of Sec.~\ref{sec:hrg}.} A sharper test would need an ensemble
in which the predicted exponent is made to move, and an analytic account of what
fixes it near $1.57$ in Eq.~\eqref{eq:core}. Neither is attempted.

\emph{Dynamics of any kind.} Everything here is a fixed point. Percolation gives
a final size and not a time course; the Ising results are equilibrium; the
epidemic chapter says so explicitly. Nothing in the book computes how long
anything takes.

\emph{Cardinality has a hard ceiling.} Equation~\eqref{eq:emitted} costs $2^{c}$
and the formalism does not degrade gracefully: it either enumerates the interior
or it does not apply. For the motifs one finds in data that is nothing, and for
a complex of thirty members it is fatal.

\section{Where this could go}
\label{sec:whereitgoes}

Four directions, of decreasing certainty.

The one the book is set up for is \emph{more models}. Chapter~\ref{ch:statmech}
takes $-\beta H$ as an argument, so a new interaction inside a complex is a
substitution and not a derivation --- which is how the unanimity rule of
Sec.~\ref{sec:unanimity} cost nothing. Anything whose interior can be enumerated
is available: Potts at integer $q$, clock models, hard-core lattice gases,
constraint satisfaction with clause structure. The cost is $2^{c}$ per complex
and $L\times L$ linear algebra thereafter, and neither depends on the rest of
the chygraph.

The one with the clearest empirical pull is \emph{chygraphs from data}.
Chapter~\ref{ch:data} showed that the objects are already there --- papers with
several authors, protein complexes, reactions, schedules --- and that a great
deal of what is recorded as a graph is the shadow of a group interaction rather
than the interaction itself. What is missing is not a method but a habit:
measuring the chy-degree ensemble, the cardinality distribution and the overlap
profile of a dataset, and reporting them, the way degree distributions are
reported now. Chapter~\ref{ch:data}'s Table~\ref{tab:real} is a start on ten
networks and no more than a start.

The one that would matter most for the theory is \emph{the hierarchy}.
Section~\ref{sec:hierarchy} argued that clustering is not one obstruction but a
sequence of them: a graph cannot see a triangle, a chygraph cannot see two
triangles sharing an edge, a region graph cannot see whatever loops the regions
form. Whether that sequence terminates for a given ensemble --- whether some
finite level captures essentially all of the correlation --- is a well-posed
question about a random geometric object, and Sec.~\ref{sec:sixty} reduces it to
the distribution of chordality-violating motifs. Answering it would say whether
the construction in this book is a rung on a ladder or the last rung that
matters.

And the most speculative is \emph{inference on the same object}. Everything here
runs the recursion forward: given an ensemble, compute a threshold. The same
machinery run backwards is what community detection and inference on networks
already do with belief propagation, and the chygraph is a strictly larger class
of models to run it on. Whether the higher-order structure buys anything there
is not obvious, and it is not obvious that it does not.

\section{The bargain, restated}
\label{sec:bargain}

The preface put it as a trade: give up the graph as the fundamental object, keep
the tree, and the apparatus comes back. That was accurate, and one qualification
should be attached to it now rather than left implicit.

What comes back is not everything, and what is given up is not only the graph.
The tree is kept by moving the loops inside complexes, where they are summed
exactly --- and complexes that share more than one node put the loops back,
one level up. On the networks that motivated the whole construction they do
share more than one node. So the chygraph moves the boundary of what counts as
local, from ``no loops'' to ``no loops between
complexes'', and buys a great deal by moving it: results that a degree-based
theory gets qualitatively wrong, on ensembles it cannot represent at all.

It does not abolish the boundary. Chapter~\ref{ch:overlap} measures what is left
on the far side of it, identifies the cause with a control that holds the
ensemble fixed, and shows what recovers the difference on one network at a time.
The ensemble version of that is the book's standing debt, and it is the right
place to stop.
